% ==============================================================================
% The Grand Daily Tournament — 2026-09-15 Match (Week 3 Day 2)
% Locked template: EB Garamond + Garamond-Math + STKaiti (v7.0 Deluxe - Unlimited Expansion Edition)
% Feature: Removed 6-page limit, expanded problem sets, heavy artillery practice
% ==============================================================================
\documentclass[a4paper,11pt]{article}

% ——— Engine & Font Setup ———
\usepackage{fontspec}
\usepackage{amsmath}
\usepackage{unicode-math}
\defaultfontfeatures{Ligatures=TeX}
\setmainfont{EBGaramond}[
  Path           = {c:/texlive/2024/texmf-dist/fonts/opentype/public/ebgaramond/},
  Extension      = .otf,
  UprightFont    = *-Regular,
  ItalicFont     = *-Italic,
  BoldFont       = *-SemiBold,
  BoldItalicFont = *-SemiBoldItalic,
  Ligatures      = TeX,
  Numbers        = OldStyle
]
\setsansfont{LibertinusSans}[
  Path        = {c:/texlive/2024/texmf-dist/fonts/opentype/public/libertinus-fonts/},
  Extension   = .otf,
  UprightFont = *-Regular,
  ItalicFont  = *-Italic,
  BoldFont    = *-Bold,
  Scale       = MatchLowercase
]
\setmathfont{Garamond-Math.otf}[
  Path  = {c:/texlive/2024/texmf-dist/fonts/opentype/public/garamond-math/},
  Scale = MatchLowercase
]

\usepackage{xeCJK}
\xeCJKsetup{
  CJKmath      = false,
  PunctStyle   = kaiming,
  CheckSingle  = true,
  AutoFakeBold = true,
  AutoFakeSlant= false
}
\setCJKmainfont{STKaiti}[
  Path        = {C:/Windows/Fonts/},
  UprightFont = STKAITI.TTF,
  BoldFont    = STKAITI.TTF,
  ItalicFont  = STKAITI.TTF,
  Script      = Default
]
\setCJKsansfont{Microsoft YaHei}
\setCJKmonofont{FangSong}

\usepackage[margin=1.4cm,headheight=14pt,headsep=10pt,footskip=18pt]{geometry}
\usepackage{booktabs,enumitem,xcolor,graphicx,tabularx}

\usepackage[breakable,skins,hooks]{tcolorbox}
\usepackage{tikz}
\usetikzlibrary{calc,arrows.meta}
\usepackage{fancyhdr,lastpage}
\usepackage{microtype}

% ——— TeXtured Colour Palette ———
\definecolor{navy}{HTML}{0F172A}
\definecolor{slate}{HTML}{1E293B}
\definecolor{royal}{HTML}{1E40AF}
\definecolor{mist}{HTML}{64748B}
\definecolor{border}{HTML}{E2E8F0}
\definecolor{paper}{HTML}{F8FAFC}
\definecolor{forest}{HTML}{065F46}
\definecolor{amber}{HTML}{D97706}
\definecolor{wine}{HTML}{991B1B}
\definecolor{ice}{HTML}{EFF6FF}
\definecolor{purple}{HTML}{6B21A8}
\definecolor{teal}{HTML}{0F766E}

% ——— Typography Tuning ———
\linespread{1.08}
\setlength{\parindent}{0pt}
\setlength{\parskip}{3pt plus 1pt minus 1pt}
\setlength{\abovedisplayskip}{4pt plus 1pt minus 1pt}
\setlength{\belowdisplayskip}{4pt plus 1pt minus 1pt}

% ——— Header / Footer ———
\pagestyle{fancy}\fancyhf{}
\fancyhead[L]{\textcolor{mist}{\footnotesize\sffamily 2027 Theoretical Physics \textperiodcentered\ 604 Calculus \& 811 Quantum Mechanics}}
\fancyhead[R]{\textcolor{slate}{\footnotesize\sffamily\bfseries Daily Tournament \textperiodcentered\ 2026-09-15}}
\fancyfoot[L]{\textcolor{mist}{\footnotesize\itshape ``Heavy Artillery Expansion \textperiodcentered\ More problems, deeper understanding.''}}
\fancyfoot[R]{\textcolor{mist}{\footnotesize\sffamily Page \thepage\ of \pageref{LastPage}}}
\renewcommand{\headrulewidth}{0.4pt}
\renewcommand{\headrule}{\color{border}\hrule width\headwidth height\headrulewidth \vskip-\headrulewidth}

% ——— tcolorbox Environments ———
\newtcolorbox{briefing}[1]{%
  enhanced, blanker, breakable,
  left=4mm, right=3mm, top=3mm, bottom=3mm,
  borderline west={2.5pt}{0pt}{royal},
  colback=paper,
  fonttitle=\sffamily\bfseries\color{slate},
  title={#1},
  attach title to upper={\par\vspace{2mm}}
}

\newtcolorbox{problembox}[3]{%
  enhanced, breakable,
  colback=white, colframe=border, boxrule=0.6pt, arc=2mm,
  left=4mm, right=4mm, top=2.5mm, bottom=2.5mm,
  fonttitle=\sffamily\bfseries\color{slate},
  title={\raisebox{0pt}{\color{#3}\rule{3.5pt}{10pt}}\hspace{4.5pt}#1\hfill{\normalfont\footnotesize\color{mist}\itshape #2}},
  colbacktitle=paper, coltitle=slate,
  titlerule=0pt, toptitle=1.8mm, bottomtitle=1.8mm
}

\newtcolorbox{solution}[1]{%
  enhanced, breakable,
  colback=white, colframe=border, boxrule=0.6pt, arc=2mm,
  left=4mm, right=4mm, top=2.2mm, bottom=2.2mm,
  fonttitle=\sffamily\bfseries\color{forest},
  title={\raisebox{0pt}{\color{forest}\rule{3pt}{10pt}}\hspace{4pt}#1},
  colbacktitle=paper, coltitle=forest,
  titlerule=0pt, toptitle=1.8mm, bottomtitle=1.8mm
}

\newtcolorbox{debrief}[1]{%
  enhanced, breakable,
  colback=white, colframe=border, boxrule=0.6pt, arc=2mm,
  left=3.5mm, right=3.5mm, top=2.5mm, bottom=2.5mm,
  fonttitle=\sffamily\bfseries\color{slate},
  title={\raisebox{0pt}{\color{amber}\rule{3pt}{10pt}}\hspace{4pt}#1},
  colbacktitle=paper, coltitle=slate,
  titlerule=0pt, toptitle=2mm, bottomtitle=2mm
}

\newcommand{\checkbox}{\ensuremath{\mdlgwhtsquare}}
\newcommand{\alerttag}[1]{\textcolor{wine}{\sffamily\bfseries [#1]}}

\begin{document}

% ==============================================================================
% PAGE 1: COVER
% ==============================================================================
\begin{titlepage}
\thispagestyle{empty}
\begin{tikzpicture}[remember picture,overlay]
  \fill[paper] (current page.south west) rectangle (current page.north east);
  \fill[navy] (current page.north west) rectangle ($(current page.north east)+(0,-3.4cm)$);
  \fill[royal] ($(current page.north west)+(0,-3.4cm)$) rectangle ($(current page.north east)+(0,-3.55cm)$);
  \fill[amber] ($(current page.north west)+(0,-3.55cm)$) rectangle ($(current page.north east)+(0,-3.62cm)$);

  \node[anchor=west, text=white] at ($(current page.north west)+(2.2cm,-1.5cm)$) {
    {\fontsize{12}{14}\selectfont\sffamily\bfseries THE GRAND DAILY TOURNAMENT \textperiodcentered\ 2027}
  };
  \node[anchor=west, text=white!50] at ($(current page.north west)+(2.2cm,-2.3cm)$) {
    {\fontsize{8.5}{11}\selectfont\sffamily UCAS \textperiodcentered\ HIAIS \textperiodcentered\ THEORETICAL PHYSICS ADVANCED TRAINING DOSSIER}
  };

  \begin{scope}[shift={($(current page.center)+(0,1.2cm)$)}]
    \draw[line width=0.5pt, border!80, dashed] (0,0) circle (3.3cm);
    \draw[line width=0.7pt, border] (0,0) circle (2.5cm);
    \draw[line width=0.4pt, royal!30] (0,0) circle (1.6cm);
    \draw[line width=1.2pt, slate!60] (-3.8,-2.2) -- (3.8,2.2);
    \draw[line width=1.2pt, slate!60] (-3.8,2.2) -- (3.8,-2.2);
    \fill[navy] (0,0) circle (3.5pt);
    \fill[amber] (0,0) circle (2pt);
  \end{scope}

  \node[anchor=center] at ($(current page.center)+(0,-3.4cm)$) {
    \begin{minipage}{0.82\textwidth}
      \centering
      {\fontsize{32}{38}\selectfont\bfseries\color{navy} Daily Tournament}\\[5pt]
      {\fontsize{12}{15}\selectfont\sffamily\bfseries\color{royal}%
        GRAND ARENA OF THEORETICAL PHYSICS \textperiodcentered\ v7.0 (Unlimited Edition)}\\[12pt]
      {\color{navy}\rule{0.5\textwidth}{1.2pt}}\\[10pt]
      {\fontsize{10.5}{14}\selectfont\color{slate}\itshape
        ``Today's contestants: Uniform Continuity, Differential Operators, Hermitian Proofs \& Past Exams''\\[3pt]
        {\small\color{mist}\upshape Closed-book \textperiodcentered\ Heavy Artillery Volume \textperiodcentered\ True Exam Simulation}}
    \end{minipage}
  };

  \node[anchor=south] at ($(current page.south)+(0,1.6cm)$) {
    \begin{minipage}{0.84\textwidth}
      \begin{tcolorbox}[
        enhanced, colback=white, colframe=border, boxrule=0.6pt, arc=2mm,
        left=4mm, right=4mm, top=3.5mm, bottom=3.5mm
      ]
        \renewcommand{\arraystretch}{1.25}
        \sffamily\small
        \begin{tabularx}{\linewidth}{@{}p{2.6cm}X@{}}
          \textbf{Date} & 2026-09-15 \quad\textbar\quad Week 3 (Day 2) \quad\textbar\quad Level 4: Angular Momentum Master\\
          \textbf{Modules} & \textbf{604}: Uniform Continuity, Improper Integrals, ODE Operators, Linear Algebra \quad\textbar\quad \textbf{811}: Hermitian Operators, Integration by parts, Real Exam \\
          \textbf{Architecture} & v7.0 Unlimited: Expansion of Volume, Introduction of Real Exam Questions\\
        \end{tabularx}
      \end{tcolorbox}
    \end{minipage}
  };
\end{tikzpicture}
\end{titlepage}

% ==============================================================================
% PAGE 2: §1 TACTICAL BRIEFING & WARM-UP RETRIEVAL
% ==============================================================================
\clearpage

\begin{center}
  {\LARGE\bfseries\color{navy} Theoretical Physics Training \textperiodcentered\ Daily Tournament}\\[3pt]
  {\small\sffamily\color{mist} 2026-09-15 \quad\textbar\quad Week 3 (Day 2) \quad\textbar\quad Phase 2: Ground-Level Foundation Reconstruction}\\[-4pt]
  {\color{navy}\rule{\textwidth}{0.8pt}}\\[-11pt]
  {\color{border}\rule{\textwidth}{0.3pt}}
\end{center}

\vspace{-6pt}

\begin{briefing}{\S 1\quad Tactical Briefing \& Warm-Up Retrieval}
\begin{itemize}[leftmargin=1.5em, itemsep=2pt, topsep=1pt]
  \item \textbf{扩容说明}：从今日起，大赛不再受6页限制，题目容量大幅提升。引入数据变换与举一反三的计算题，以及量子力学真题模拟，要求强化运算肌肉记忆。
  \item \textbf{604 今日靶向}：一致连续性构造与反例判定；反常积分审敛法的实际应用；可降阶微分方程求解；微分算子法求解二阶常系数线性非齐次微分方程。
  \item \textbf{811 今日靶向}：厄米算符数学定义、利用无穷远边界条件进行分部积分的严谨推导；动量算符的厄米性证明；【新增】真题级别算符对易子实战演练。
\end{itemize}

\vspace{1mm}
\textbf{Warm-Up Retrieval Checklist} (5--8\,min): without notes, write on scratch paper:\\
\checkbox\ \alerttag{604} 微分算子法公式：$\frac{1}{f(D)} e^{\lambda x} = \frac{1}{f(\lambda)} e^{\lambda x}$ （当 $f(\lambda) \neq 0$）\quad
\checkbox\ \alerttag{604} p级数反常积分收敛条件：$\int_1^{\infty} \frac{1}{x^p} dx$ 收敛当且仅当 $p > 1$ \quad
\checkbox\ \alerttag{811} 厄米算符定义：$\int \psi_1^* (\hat{F}\psi_2) d\tau = \int (\hat{F}\psi_1)^* \psi_2 d\tau$ \quad
\checkbox\ \alerttag{811} 动量算符：$\hat{p} = -i\hbar \nabla$
\end{briefing}

\vspace{4pt}

\begin{tcolorbox}[
  enhanced, colback=white, colframe=border, boxrule=0.6pt, arc=2mm,
  left=4mm, right=4mm, top=3mm, bottom=3mm,
  fonttitle=\sffamily\bfseries\color{slate},
  title={\raisebox{0pt}{\color{royal}\rule{3pt}{10pt}}\hspace{4pt}Warm-Up Retrieval Workspace \textbar\ 白纸破冰与间隔提取草稿区},
  colbacktitle=paper, coltitle=slate,
  titlerule=0pt, toptitle=2mm, bottomtitle=2mm
]
\textcolor{mist}{\footnotesize\itshape 5--8 分钟闭卷盲写：默写上述公式，预热大脑。}
\vspace{8.0cm}
\end{tcolorbox}

% ==============================================================================
% PAGE 3: §2 SECTION A & SECTION B (CALCULUS & ODE)
% ==============================================================================
\clearpage

{\Large\sffamily\bfseries\color{navy} \S 2\quad Foundation Matrix \& Concept Blitz}
\hfill{\small\sffamily\color{mist}\itshape Complete independently under timed conditions}

\vspace{6pt}

{\normalsize\sffamily\bfseries\color{royal} \S A\quad 高数基础区 (Calculus Core \textperiodcentered\ 30 min)}
\vspace{2pt}

\begin{problembox}{Problem A.1 \textbar\ 604 Calculus: 一致连续性判别}{Suggested time: 10\,min}{royal}
判断函数 $f(x) = x \sin\left(\frac{1}{x}\right)$ 在区间 $(0, 1]$ 上是否一致连续？\\
判断函数 $g(x) = x^2$ 在区间 $(0, \infty)$ 上是否一致连续？若不是，请构造反例证明。
\end{problembox}
\vspace{6.5cm}

\begin{problembox}{Problem A.2 \textbar\ 604 Calculus: 反常积分审敛}{Suggested time: 10\,min}{royal}
判别反常积分的收敛性：
\[
  I_1 = \int_1^{\infty} \frac{\ln x}{1+x^2} dx, \quad I_2 = \int_0^1 \frac{1}{\sqrt{x(1-x)}} dx
\]
\end{problembox}
\vspace{6.5cm}

\clearpage

{\normalsize\sffamily\bfseries\color{amber} \S B\quad 微分方程与线代进阶区 (ODE \& LA Blitz \textperiodcentered\ 30 min)}
\vspace{2pt}

\begin{problembox}{Problem B.1 \textbar\ 604 ODE: 缺项可降阶方程}{Suggested time: 12\,min}{amber}
求解以下不显含 $y$ 的可降阶非线性微分方程，并求其满足初值条件 $y(0)=1, y'(0)=1$ 的特解：
\[
  y'' + (y')^2 = 1
\]
\textit{提示：令 $p = y'$，转化为关于 $p(x)$ 的一阶方程，注意分离变量。}
\end{problembox}
\vspace{8.5cm}

\begin{problembox}{Problem B.2 \textbar\ 604 ODE: 微分算子法实战}{Suggested time: 12\,min}{amber}
利用微分算子法（Differential Operator Method），求二阶常系数线性非齐次微分方程的通解：
\[
  y'' - 4y' + 4y = 8e^{2x} + 4x^2
\]
\textit{提示：右端为指数项与多项式项的叠加，利用共振公式 $\frac{1}{(D-a)^n} e^{ax} = \frac{x^n}{n!} e^{ax}$ 以及泰勒展开。}
\end{problembox}
\vspace{8.5cm}

\clearpage
\begin{problembox}{Problem B.3 \textbar\ 604 Linear Algebra: 矩阵运算与特征值}{Suggested time: 10\,min}{amber}
设矩阵 $A$ 满足 $A^2 - 3A + 2E = 0$（其中 $E$ 为单位矩阵），且 $A$ 是 $3 \times 3$ 的实矩阵。
已知 $A$ 不等于 $E$ 也不等于 $2E$，求 $A$ 可能的特征值，并列举出 $A$ 迹 (Trace) 的所有可能值。
\end{problembox}
\vspace{10.0cm}

% ==============================================================================
% PAGE : §3 SECTION C (QUANTUM MECHANICS)
% ==============================================================================
\clearpage

{\Large\sffamily\bfseries\color{navy} \S 3\quad Quantum Mechanics Core \& Real Exams}
\hfill{\small\sffamily\color{mist}\itshape Heavy Artillery Volume}

\vspace{6pt}
{\normalsize\sffamily\bfseries\color{forest} \S C\quad 厄米算符与对易子实战 (Hermitian \& Commutator \textperiodcentered\ 30 min)}
\vspace{2pt}

\begin{problembox}{Problem C.1 \textbar\ 811 QM: 厄米算符的严谨证明}{Suggested time: 15\,min}{forest}
在一维空间 $(-\infty, +\infty)$ 中，已知动量算符为 $\hat{p} = -i\hbar \frac{d}{dx}$。波函数 $\psi(x)$ 满足无穷远处的自然边界条件（即 $x \to \pm\infty$ 时，$\psi(x) \to 0$）。
\begin{enumerate}
    \item 严格利用分部积分法证明 $\hat{p}$ 是厄米算符。
    \item 证明若 $\hat{A}$ 和 $\hat{B}$ 均为厄米算符，则算符 $\hat{C} = i[\hat{A}, \hat{B}]$ 也是厄米算符。
\end{enumerate}
\end{problembox}
\vspace{10.0cm}

\clearpage
\begin{problembox}{Problem C.2 \textbar\ 811 QM: 真题模拟训练 (Adapted from UCAS Past Exams)}{Suggested time: 15\,min}{forest}
（真题风格演练）在量子力学中，算符之间的对易关系对于确定不确定性原理具有核心意义。
已知坐标算符 $\hat{x}$ 与动量算符 $\hat{p}$ 满足基本对易关系 $[\hat{x}, \hat{p}] = i\hbar$。
\begin{enumerate}
    \item 求算符对易子 $[\hat{x}, \hat{p}^2]$ 的表达式。
    \item 推导 $[\hat{x}^2, \hat{p}^2]$ 的结果，将其表示为 $\hat{x}$ 与 $\hat{p}$ 的函数。
    \item 结合广义不确定性原理关系式 $\Delta A \Delta B \ge \frac{1}{2}|\langle [\hat{A}, \hat{B}] \rangle|$，简述第一问结果的物理意义。
\end{enumerate}
\end{problembox}
\vspace{12.0cm}

% ==============================================================================
% PAGE : DEBRIEFING AREA
% ==============================================================================
\clearpage

{\Large\sffamily\bfseries\color{navy} \S 4\quad Post-Match Debriefing \& Error Analysis}

\vspace{4pt}

\begin{debrief}{Debriefing Protocol \textbar\ 复盘与反思 (Fill after grading)}
\textbf{1. 知识盲区与计算失误 (Blind Spots \& Execution Errors):}
\vspace{3.5cm}

\textbf{2. 概念纠偏 (Conceptual Corrections):}
\vspace{3.5cm}

\textbf{3. 下一步行动计划 (Next-Step Action Items):}
\vspace{3.5cm}
\end{debrief}

\vspace{4pt}
\begin{center}
\sffamily\color{mist}\small
End of Tournament. Stand tall, review rigorously, and prepare for the next battle.\\
\textit{May the Hamiltonian guide your path.}
\end{center}



% ==============================================================================
% PAGE : SOLUTIONS
% ==============================================================================
\clearpage
{\Large\sffamily\bfseries\color{navy} \S 5\quad 官方解答与解析 (Official Solutions)}
\vspace{6pt}

\begin{solution}{A.1 \textbar\ 604 一致连续性判别}
1. 对于 $f(x) = x \sin(1/x)$ 在 $(0, 1]$ 上：\\
由于 $\lim_{x \to 0} x \sin(1/x) = 0$ 存在且有限，我们可以补充定义 $f(0)=0$，使其在闭区间 $[0,1]$ 上连续。由 Cantor 定理（闭区间上的连续函数必一致连续）可知，扩展后的函数在 $[0,1]$ 上一致连续。因此它在子区间 $(0,1]$ 上也一致连续。

2. 对于 $g(x) = x^2$ 在 $(0, \infty)$ 上：\\
不一致连续。构造反例序列：取 $x_n = \sqrt{n+1}$， $y_n = \sqrt{n}$。\\
当 $n \to \infty$ 时，$|x_n - y_n| = \frac{1}{\sqrt{n+1}+\sqrt{n}} \to 0$。\\
但对应的函数值之差为 $|g(x_n) - g(y_n)| = |(n+1) - n| = 1$。\\
自变量距离趋于 0，而函数值距离无法趋于 0（恒为 1），违背了一致连续性的定义 $\forall \epsilon, \exists \delta$，故不一致连续。
\end{solution}

\begin{solution}{A.2 \textbar\ 604 反常积分审敛}
1. $I_1 = \int_1^{\infty} \frac{\ln x}{1+x^2} dx$：\\
被积函数在 $x \to \infty$ 时，渐近行为是 $\sim \frac{\ln x}{x^2}$。\\
选取比较函数 $g(x) = \frac{1}{x^{3/2}}$，利用极限比较法：
$\lim_{x \to \infty} \frac{(\ln x)/(1+x^2)}{1/x^{3/2}} = \lim_{x \to \infty} \frac{x^{3/2}\ln x}{x^2} = \lim_{x \to \infty} \frac{\ln x}{x^{1/2}} = 0$。\\
因为 $p=3/2 > 1$，积分 $\int_1^\infty \frac{1}{x^{3/2}} dx$ 收敛，故原积分收敛。

2. $I_2 = \int_0^1 \frac{1}{\sqrt{x(1-x)}} dx$：\\
积分在 $x=0$ 和 $x=1$ 处均为瑕点，拆分为 $\int_0^{1/2} + \int_{1/2}^1$。\\
在 $x \to 0^+$ 时，$\frac{1}{\sqrt{x(1-x)}} \sim \frac{1}{x^{1/2}}$，$p=1/2 < 1$，瑕积分收敛。\\
在 $x \to 1^-$ 时，令 $t=1-x$，$t \to 0^+$，被积函数 $\sim \frac{1}{t^{1/2}}$，同理 $p=1/2 < 1$ 收敛。\\
两部分均收敛，故该反常积分收敛。
\end{solution}

\begin{solution}{B.1 \textbar\ 604 可降阶微分方程}
令 $p = y'$，原方程化为一阶方程：$p' + p^2 = 1$。\\
注意初值条件：$y'(0)=1 \implies p(0)=1$。\\
将 $p(0)=1$ 代入方程 $p' = 1 - p^2$，得到 $p'(0) = 0$。\\
如果进行变量分离 $\frac{dp}{1-p^2} = dx$，由于 $1-p^2$ 为 0，分母为 0 无法直接积分！这是一个极其容易踩坑的陷阱点。\\
事实上，显然 $p(x) \equiv 1$ 是该微分方程的常数解（平衡解），且满足初值条件。\\
因此 $y' = 1 \implies y = x + C$。\\
代入 $y(0)=1$ 得 $C=1$，故特解为 $y = x + 1$。
\end{solution}

\begin{solution}{B.2 \textbar\ 604 微分算子法}
特征方程 $r^2 - 4r + 4 = 0$，二重根 $r=2$，齐次通解 $y_h = (C_1 + C_2 x)e^{2x}$。\\
特解 $y_{p1} = \frac{1}{(D-2)^2} 8e^{2x}$，由于 $\lambda=2$ 是特征方程的二重根（发生共振），运用位移公式：\\
$y_{p1} = 8 \frac{x^2}{2!} e^{2x} = 4x^2 e^{2x}$。\\
特解 $y_{p2} = \frac{1}{D^2 - 4D + 4} 4x^2 = \frac{1}{4(1 - D/2)^2} 4x^2 = (1 - D/2)^{-2} x^2$。\\
由泰勒展开 $(1-X)^{-2} = 1 + 2X + 3X^2 + O(X^3)$，代入 $X = D/2$：\\
$y_{p2} = (1 + 2(D/2) + 3(D/2)^2) x^2 = (1 + D + \frac{3}{4}D^2) x^2 = x^2 + 2x + \frac{3}{2}$。\\
故通解为 $y = (C_1 + C_2 x)e^{2x} + 4x^2 e^{2x} + x^2 + 2x + \frac{3}{2}$。
\end{solution}

\begin{solution}{B.3 \textbar\ 604 矩阵迹与特征值}
由 $A^2 - 3A + 2E = 0$ 得 $(A-E)(A-2E)=0$。\\
矩阵 $A$ 的极小多项式必须整除 $\lambda^2 - 3\lambda + 2 = (\lambda-1)(\lambda-2)$，故特征值只能取 1 或 2。\\
由于 $A$ 是 $3 \times 3$ 矩阵，共有三个特征值，其可能的组合（考虑重数）有：\\
\{1,1,1\}, \{1,1,2\}, \{1,2,2\}, \{2,2,2\}。\\
若特征值为 \{1,1,1\}，因极小多项式无重根，矩阵必可对角化，得 $A = P E P^{-1} = E$，与已知 $A \neq E$ 矛盾。\\
若特征值为 \{2,2,2\}，同理 $A = 2E$，与已知矛盾。\\
故特征值的组合只能是 \{1,1,2\} 或 \{1,2,2\}。\\
矩阵的迹等于特征值之和：\\
若为 \{1,1,2\}，则 $Tr(A) = 1+1+2 = 4$。\\
若为 \{1,2,2\}，则 $Tr(A) = 1+2+2 = 5$。\\
所以 $A$ 的迹的所有可能值为 4 和 5。
\end{solution}

\begin{solution}{C.1 \textbar\ 811 厄米算符证明}
1. $\hat{p}$ 的厄米性：\\
$\int_{-\infty}^{\infty} \psi_1^* (\hat{p} \psi_2) dx = \int_{-\infty}^{\infty} \psi_1^* \left(-i\hbar \frac{d\psi_2}{dx}\right) dx$\\
利用分部积分：
$= \left[-i\hbar \psi_1^* \psi_2 \right]_{-\infty}^{\infty} - \int_{-\infty}^{\infty} \left(-i\hbar \frac{d\psi_1^*}{dx}\right) \psi_2 dx$\\
因为波函数在无穷远处趋于 0，第一项边界项为 0。\\
剩下的积分为：$\int_{-\infty}^{\infty} \left(i\hbar \frac{d\psi_1}{dx}\right)^* \psi_2 dx = \int_{-\infty}^{\infty} (\hat{p} \psi_1)^* \psi_2 dx$。\\
满足定义 $\int \psi_1^* (\hat{p} \psi_2) = \int (\hat{p} \psi_1)^* \psi_2$，故 $\hat{p}$ 为厄米算符。

2. $\hat{C} = i[\hat{A}, \hat{B}]$ 的厄米性：\\
求 $\hat{C}$ 的厄米共轭：
$\hat{C}^\dagger = \left( i(\hat{A}\hat{B} - \hat{B}\hat{A}) \right)^\dagger = -i (\hat{A}\hat{B} - \hat{B}\hat{A})^\dagger$\\
利用性质 $(\hat{X}\hat{Y})^\dagger = \hat{Y}^\dagger \hat{X}^\dagger$ 和 $\hat{A}^\dagger = \hat{A}$：\\
$= -i ( (\hat{A}\hat{B})^\dagger - (\hat{B}\hat{A})^\dagger ) = -i ( \hat{B}^\dagger \hat{A}^\dagger - \hat{A}^\dagger \hat{B}^\dagger ) = -i (\hat{B}\hat{A} - \hat{A}\hat{B})$\\
$= i(\hat{A}\hat{B} - \hat{B}\hat{A}) = \hat{C}$。\\
所以 $\hat{C}^\dagger = \hat{C}$，即 $\hat{C}$ 也是厄米算符。
\end{solution}

\begin{solution}{C.2 \textbar\ 811 算符对易子}
已知 $[\hat{x}, \hat{p}] = i\hbar$。\\
1. $[\hat{x}, \hat{p}^2] = \hat{p}[\hat{x}, \hat{p}] + [\hat{x}, \hat{p}]\hat{p} = \hat{p}(i\hbar) + (i\hbar)\hat{p} = 2i\hbar\hat{p}$。\\
2. $[\hat{x}^2, \hat{p}^2] = \hat{x}[\hat{x}, \hat{p}^2] + [\hat{x}, \hat{p}^2]\hat{x}$\\
将第一问的结果代入：$= \hat{x}(2i\hbar\hat{p}) + (2i\hbar\hat{p})\hat{x} = 2i\hbar(\hat{x}\hat{p} + \hat{p}\hat{x})$。\\
3. 物理意义：由广义不确定性关系 $\Delta x \Delta (p^2) \ge \frac{1}{2} | \langle [\hat{x}, \hat{p}^2] \rangle | = \hbar |\langle \hat{p} \rangle|$。\\
这表明，粒子的动量期望值绝对值越大，其位置和动能（正比于 $p^2$）的联合不确定度下限就越大。换言之，高速运动粒子的位置和动能更难以被同时精确测量。
\end{solution}

\end{document}
